\chapter{Introducción}
\label{chap:intro}

\drop{L}{a} calidad de los datos constituye uno de los pilares fundamentales sobre
los que se sustenta el éxito de cualquier proyecto de \ac{IA} o \ac{ML}.
A medida que las organizaciones adoptan sistemas basados en datos para la toma
de decisiones estratégicas, la necesidad de garantizar que dichos datos sean
completos, precisos, consistentes y oportunos se ha convertido en un requisito
ineludible.  Sin embargo, la realidad de la industria muestra que la calidad de
los datos sigue siendo un problema generalizado y costoso, tanto en términos
económicos como en sus consecuencias sobre la fiabilidad de los modelos
resultantes.  El presente \ac{TFG} aborda esta problemática mediante el diseño
e implementación de \dataqual, una plataforma web integral para la evaluación
automatizada de la calidad de datos en proyectos de \ac{IA}.

Este capítulo introductorio ofrece una visión general del problema que motiva
el desarrollo de la plataforma, presenta la solución propuesta y sus
principales capacidades, justifica la necesidad de una herramienta dedicada
frente a los enfoques tradicionales, identifica el público objetivo al que se
dirige y describe la estructura del resto del documento.


%% ---------------------------------------------------------------------------
\section{Motivación}
\label{sec:motivacion}

El principio conocido como \emph{garbage in, garbage out} resume de forma
elocuente una verdad ampliamente reconocida en el ámbito de la ciencia de
datos: la calidad de las salidas de un sistema está intrínsecamente limitada
por la calidad de sus entradas.  En el contexto de la \ac{IA} y el \ac{ML},
este principio adquiere una relevancia especial, ya que los modelos aprenden
patrones directamente de los datos de entrenamiento~\citep{wang1996}.  Si
estos datos contienen errores, sesgos o valores ausentes, los modelos
resultantes heredarán, y con frecuencia amplificarán, dichas
deficiencias~\citep{redman1998}.

Diversos estudios han demostrado que los problemas de calidad de datos
constituyen la causa principal de fallos de sistemas de \ac{IA} en
producción, por encima incluso de los errores algorítmicos.
Sambasivan et~al.~(\citeyear{sambasivan2021}) documentaron, a partir de entrevistas
con 53 profesionales de \ac{IA} en múltiples organizaciones, que las
deficiencias en los datos (etiquetado incorrecto, datos
desactualizados, muestras no representativas) son la causa
predominante de los fracasos en despliegues de \ac{ML}.  A nivel técnico, la presencia de etiquetas incorrectas en los datos de
entrenamiento puede degradar significativamente la precisión del modelo
resultante, una degradación que ninguna mejora arquitectónica puede compensar
de forma fiable~\citep{ng2021datacentric}.

Las consecuencias económicas de esta situación son igualmente significativas:
el coste de corregir errores crece exponencialmente a medida que avanzan por
las etapas de un flujo de procesamiento~\citep{redman2001}.  Encuestas a profesionales del sector revelan consistentemente que una
proporción significativa de su tiempo se dedica a tareas de limpieza y
preparación de datos~\citep{kaggle2022}, en detrimento del tiempo disponible
para el análisis y el desarrollo de modelos, lo que evidencia la necesidad de
herramientas que automaticen la evaluación de la calidad.

La calidad de los datos no es, además, únicamente una cuestión técnica.
Los datos desequilibrados o no representativos son una fuente frecuente de
sesgo con consecuencias éticas y sociales: el algoritmo de selección de
personal de Amazon, descartado en 2018 por discriminar sistemáticamente a las
candidatas femeninas~\citep{amazon2018}, y los sistemas de reconocimiento
facial con tasas de error desproporcionadas para personas de piel oscura
documentados por Buolamwini y Gebru~(\citeyear{buolamwini2018}) ilustran cómo las
deficiencias en los datos de entrenamiento se trasladan al comportamiento del
modelo.  El marco regulatorio refuerza esta dimensión: el Reglamento Europeo
de \ac{IA}~\citep{aiact2024} establece requisitos explícitos de calidad,
trazabilidad y documentación de los datos en sistemas de alto riesgo, y el
\ac{RGPD}~\citep{rgpd2016} exige que los datos personales sean exactos y estén
actualizados, convirtiendo la evaluación formal de la calidad en una
obligación legal, no solo una buena práctica.

Este trabajo nace del interés personal del autor por el tratamiento y la
gestión de los datos como disciplina, y de su voluntad de profundizar en el
ámbito de la inteligencia artificial y la ciencia de datos, campos en los que
la calidad del dato de entrada resulta, precisamente, el factor más
determinante.


%% ---------------------------------------------------------------------------
\section{Presentación de la solución}
\label{sec:solucion}

Para abordar la problemática descrita, se ha diseñado e implementado
\dataqual, una plataforma web \emph{full-stack} para la evaluación
automatizada de la calidad de datos en proyectos de \ac{IA}.  A diferencia de las herramientas existentes, \dataqual
integra en un único entorno la configuración de métricas, la ejecución de
evaluaciones y la visualización de resultados, sin necesidad de escribir
código.  La plataforma combina una interfaz gráfica interactiva con
una \ac{API} \ac{REST}ful~\citep{fielding2000} que permite su integración en
flujos de trabajo existentes, y se despliega como ocho servicios Docker
orquestados con Docker~Compose para garantizar portabilidad y soberanía
sobre los datos.  Sus capacidades principales son:

\begin{itemize}
  \item \textbf{Carga de conjuntos de datos} en formato \ac{CSV}, con
    almacenamiento seguro en MinIO y versionado gestionado mediante
    una jerarquía de versiones con etiquetas personalizables.

  \item \textbf{Perfilado exploratorio del dataset}: análisis estadístico
    automático por columna que incluye distribuciones, valores nulos,
    tipos de datos, estadísticos descriptivos e histogramas, para facilitar
    la comprensión del dato antes de configurar las métricas formales.

  \item \textbf{Configuración de métricas de calidad} adaptadas al tipo de
    proyecto, con parámetros y umbrales personalizables alineados con los
    estándares \ac{ISO}/IEC~25012~\citep{iso25012} y
    \ac{ISO}/IEC~5259~\citep{iso5259}.

  \item \textbf{Ejecución asíncrona de evaluaciones}, que analiza los
    conjuntos de datos frente a las métricas configuradas sin bloquear la
    interfaz de usuario.

  \item \textbf{Visualización interactiva de resultados} mediante cuadros de
    mando que incluyen gráficos de evolución, tablas de métricas por columna,
    histogramas, diagramas de caja y matrices de correlación.

  \item \textbf{Seguimiento de la evolución de la calidad} a lo largo del
    tiempo, gracias a un sistema de \emph{fingerprinting} que permite comparar
    automáticamente los resultados con líneas base anteriores.

  \item \textbf{Quality Gates} (puertas de calidad): umbrales mínimos
    configurables que determinan si un conjunto de datos es apto para su uso,
    con estados \texttt{PASSED}, \texttt{WARNING} y \texttt{FAILED},
    inspirados en el concepto homónimo de SonarQube~\citep{sonarqube}.

  \item \textbf{Protección de datos sensibles}: el usuario marca las columnas
    que contienen información sensible al cargar el dataset, y la plataforma
    enmascara automáticamente sus valores en todos los resultados y muestras
    de la evaluación, evitando su exposición en los informes.
\end{itemize}


%% ---------------------------------------------------------------------------
\section{Justificación de la solución}
\label{sec:justificacion}

Los enfoques tradicionales para la evaluación de la calidad de datos,
\emph{scripts} ad~hoc en Python, consultas \ac{SQL} manuales o cuadernos
Jupyter, presentan limitaciones que dificultan su adopción sistemática en
entornos profesionales.  La primera de ellas es la falta de estandarización:
cada equipo define sus propias métricas y umbrales de forma aislada, lo que
impide la comparación entre proyectos y complica la auditoría.  \dataqual
aborda esto mediante un catálogo alineado con \ac{ISO}/IEC~25012~\citep{iso25012}
y la serie \ac{ISO}/IEC~5259~\citep{iso5259}, que proporciona un marco de
referencia común basado en el modelo \ac{SQuaRE}.  A esto se añade la
ausencia de historización: los \emph{scripts} evalúan instantáneas puntuales
sin conservar un registro que permita detectar degradaciones progresivas de
la calidad, algo que \dataqual resuelve mediante el seguimiento automático
entre ejecuciones.

Una segunda barrera es de naturaleza técnica.  Herramientas como Great
Expectations~\citep{greatexpectations} o Deequ de Amazon requieren
conocimientos de programación que no todos los perfiles involucrados en un
proyecto de datos poseen, y los \emph{scripts} independientes tampoco se
integran de forma natural en \emph{pipelines} existentes.  \dataqual ofrece
una interfaz web accesible para usuarios no técnicos y, simultáneamente, una
\ac{API} \ac{REST} para la integración programática con flujos de \ac{MLOps}
y procesos \ac{ETL}.

Por último, ninguno de estos enfoques proporciona una señal clara de
aprobación o rechazo sobre la calidad del conjunto de datos.  Esta carencia
es especialmente crítica en flujos de \ac{MLOps}, donde se necesita un
criterio objetivo para decidir si los datos son aptos antes de entrenar un
modelo.  Inspirándose en el concepto de \emph{Quality Gate} de
SonarQube~\citep{sonarqube}, \dataqual implementa puertas de calidad
configurables con estados \texttt{PASSED}, \texttt{WARNING} y \texttt{FAILED}
que convierten una evaluación cuantitativa en una decisión accionable.


%% ---------------------------------------------------------------------------
\section{Público objetivo}
\label{sec:publico}

\dataqual está dirigida a los distintos perfiles profesionales que
intervienen en el ciclo de vida de un proyecto de datos.  Los ingenieros de
datos pueden integrar las evaluaciones en sus \emph{pipelines} \ac{ETL}
mediante la \ac{API} \ac{REST}, mientras que los científicos de datos e
ingenieros de \ac{ML} disponen de un mecanismo sistemático para verificar la
calidad de los datos de entrenamiento antes de construir y desplegar modelos.
Los analistas sin perfil técnico disponen de una interfaz visual para
explorar la calidad de los conjuntos de datos de forma autónoma.
Finalmente, los responsables de calidad y cumplimiento normativo encuentran
en la plataforma evidencias auditables que facilitan la justificación ante
los requisitos del \ac{RGPD} y el Reglamento Europeo de \ac{IA}, y los
directores técnicos disponen de cuadros de mando agregados para la toma de
decisiones informadas.


%% ---------------------------------------------------------------------------
\section{Alcance del proyecto}
\label{sec:alcance}

\dataqual es una plataforma \emph{full-stack} orientada a proyectos de \ac{IA}
y \ac{ML} de escala mediana, con soporte para conjuntos de datos tabulares de
hasta 100\,MB en formato \ac{CSV}.  La
plataforma cubre el ciclo completo de evaluación: desde la carga y el
versionado de los datos hasta la emisión de un veredicto automatizado mediante
\emph{Quality Gates}, pasando por la ejecución de seis métricas de calidad
alineadas con \ac{ISO}/IEC~5259, el perfilado exploratorio, el seguimiento de
incidencias entre evaluaciones y la protección de datos personales.  El
sistema se despliega en entorno local mediante Docker~Compose y expone una
\ac{API} \ac{REST} con autenticación \ac{JWT} para su integración en flujos de
trabajo externos.

Quedan \textbf{fuera del alcance} del presente trabajo: el procesamiento de
conjuntos de datos superiores a 100\,MB, el análisis de flujos en tiempo real
(\emph{streaming}), la corrección automática de errores, la integración nativa
con plataformas \emph{cloud} propietarias (AWS, Azure, GCP) y la
internacionalización de la interfaz.

%% ---------------------------------------------------------------------------
\section{Estructura del documento}
\label{sec:estructura}

El presente documento se organiza en siete capítulos y siete anexos.  A
continuación se describe brevemente el contenido de cada uno de los capítulos
restantes:

\begin{definitionlist}
\item[Capítulo~\ref{chap:antecedentes}: \nameref{chap:antecedentes}]
  Revisa los fundamentos teóricos de la calidad de datos, los principales
  estándares internacionales (\ac{ISO}/IEC~25012, \ac{ISO}/IEC~5259), las
  dimensiones de calidad relevantes y las herramientas existentes en el
  mercado, identificando las carencias que justifican el desarrollo de
  \dataqual.

\item[Capítulo~\ref{chap:objetivos}: \nameref{chap:objetivos}]
  Define el objetivo general y los objetivos específicos del proyecto,
  estableciendo el alcance detallado y las restricciones del trabajo.

\item[Capítulo~\ref{chap:metodologia}: \nameref{chap:metodologia}]
  Describe la metodología de gestión y desarrollo adoptada, la pila
  tecnológica, el backlog inicial, la planificación (Gantt), la
  estimación de esfuerzo, los costes y la gestión de riesgos.

\item[Capítulo~\ref{chap:desarrollo}: \nameref{chap:desarrollo}]
  Desarrolla el proceso de implementación iterativo por sprints, incluyendo
  para cada sprint la planificación, el trabajo realizado, las decisiones
  técnicas adoptadas y la retrospectiva.

\item[Capítulo~\ref{chap:resultados}: \nameref{chap:resultados}]
  Presenta los resultados globales obtenidos, analiza el cumplimiento de los
  objetivos planteados y compara \dataqual con las herramientas existentes.

\item[Capítulo~\ref{chap:conclusiones}: \nameref{chap:conclusiones}]
  Recoge las conclusiones, la revisión de objetivos, la aportación del
  proyecto, las competencias adquiridas, la valoración personal y las líneas
  de trabajo futuro.
\end{definitionlist}

El documento incluye siete anexos: especificación completa de la \ac{API}
\ac{REST} (Anexo~A), catálogo de métricas de calidad (Anexo~B), backlog
completo (Anexo~C), plan de riesgos detallado (Anexo~D), estimación de costes
y viabilidad (Anexo~E), manual de usuario (Anexo~F) y colección de diagramas
\ac{UML} (Anexo~G).


% Local Variables:
%  coding: utf-8
%  mode: latex
%  mode: flyspell
%  ispell-local-dictionary: "castellano8"
% End:
